\begin{abstract}
Uncertainty controls determine how an agent learns inside a world model. Ending a rollout assigns a value to its boundary. Weighting an advantage changes how individual decisions contribute to the policy update. We analyze these two choices directly. For stopped trajectories, we derive the exact boundary-value error and characterize when a constant reward shift preserves policy rankings. For a finite-action decision state, we characterize when fixed positive weights produce nonnegative alignment with the return gradient for every exact centered advantage. The characterization separates compatibility with the policy's score space from the geometry used to update its parameters. Fisher preconditioning resolves the geometric obstruction when the compatibility condition holds. Explicit finite-action and Gaussian examples show how positive weights can reverse the return gradient. An existing Pendulum study provides a continuous-control setting: four imagined-refinement rules reduce mean environment return from a weaker starting policy, and uncertainty-aware rules have small return gaps alongside poor control. The analysis gives concrete conditions for choosing rollout boundaries and weighting rules while keeping the intended learning objective in view.
\end{abstract}

\section{Introduction}

A world model gives a control policy a place to practice. The policy can explore predicted trajectories and update its behavior before interacting with the environment again. This makes the rule for using those trajectories part of the learning problem. A short rollout chooses which future rewards enter a target. A confidence weight changes which decisions dominate an update. Each choice deserves an analysis of the objective it creates.

This paper asks: \emph{how do uncertainty-based stopping and advantage weighting change the quantity that policy learning improves?} We study the boundary of an imagined trajectory and the direction of a weighted policy gradient. The resulting conditions can be checked in small analytical models and mapped to specific implementation decisions.

The motivating study implements continuous-control PPO, a five-member dynamics ensemble, and four imagined-refinement rules in \code{Pendulum-v1}. Starting from a partially trained policy with environment return $-576.4$, the four rules produce lower mean returns. Uncertainty termination yields $-1472.7\pm24.1$, while weighted advantages yield $-1462.1\pm4.2$. Both have relatively small model--environment return gaps. These outcomes motivate a closer look at the learning rule, alongside prediction accuracy.

Two structural observations organize the analysis. First, ending a segment with zero continuation deletes its remaining value. In a task with nonpositive rewards, that deletion adds the expected omitted cost to the stopped return. If the stopping time depends on the policy, it can also change policy rankings. Second, positive weights rescale individual score contributions. Their sum can point against the original return gradient. The relevant condition concerns the score space and the update metric.

The paper develops three results:
\begin{enumerate}
\item \textbf{Stopping boundaries and reward shifts.} We express the entire stopped-return error through a boundary value. We characterize reward-shift invariance and give the boundary compensation that restores it on every path.
\item \textbf{A complete finite-action alignment criterion.} For a fixed full-support policy at one decision state and fixed positive weights, we give necessary and sufficient conditions for nonnegative alignment with the return gradient for every exact centered advantage. Under the compatibility condition, Fisher preconditioning gives strict ascent whenever the original gradient is nonzero.
\item \textbf{Exact examples and an implementation connection.} Finite-action examples isolate two ways the alignment criterion can fail. A bounded-reward Gaussian construction connects the direction reversal to a continuous actor. Existing Pendulum results show the behavior of the implemented stopping and weighting rules.
\end{enumerate}

\subsection{Related work}

World-model-based policy refinement has a broad control motivation. PETS uses probabilistic dynamics ensembles and trajectory sampling \cite{pets}. MBPO analyzes model use and motivates short rollouts branched from real experience \cite{mbpo}. World4RL and WoVR apply learned visual simulators to policy refinement for robotic manipulation \cite{world4rl,wovr}. Our setting makes boundary targets and update directions explicit in a small state-space control problem.

Pardo et al. distinguish a task horizon from a training truncation and analyze partial-episode bootstrapping \cite{pardo}. MOReL assigns a low-reward absorbing state to unknown regions, while MOPO uses uncertainty penalties to construct a conservative objective \cite{morel,mopo}. TATU combines uncertainty-driven trajectory truncation with penalty terms \cite{tatu}. These methods make the value assigned to uncertain experience part of their construction. Section~\ref{sec:stopping} develops an exact finite-horizon comparison and a reward-translation criterion for a specified boundary rule.

Policy-gradient theory relates expected-return derivatives to score functions \cite{sutton}; natural gradients use the Fisher geometry of the policy family \cite{kakade}. The distinction between unrestricted and restricted policy classes is central to modern analyses \cite{agarwal}. Action-dependent baseline estimators and their implementation effects are studied by Tucker et al. \cite{tucker}. Advantage-weighted regression uses a weighted fitting objective \cite{awr}. Osband analyzes positive gating, variance, and directional behavior in specific bandit and learning settings \cite{osband}. Han and Hu study reward-shift invariance and score centering under regularized importance weights \cite{han}. Section~\ref{sec:weighting} fixes the weight and policy score space, then quantifies over every centered advantage. This gives an algebraic criterion for alignment with the original return gradient.
